\documentclass[12pt]{article}
\usepackage{amsmath,amssymb,amsthm}
\usepackage{geometry}
\geometry{margin=1in}
\usepackage{hyperref}
\usepackage{booktabs}
\usepackage{authblk}

\newtheorem{conjecture}{Conjecture}
\newtheorem{remark}[conjecture]{Remark}
\newtheorem{corollary}[conjecture]{Corollary}

\newcommand{\Z}{\mathbb{Z}}
\newcommand{\R}{\mathbb{R}}
\newcommand{\C}{\mathbb{C}}
\newcommand{\ph}{\varphi}
\newcommand{\Lk}{L_k}
\newcommand{\logph}{\log\ph}

\title{Hyperbolic Flavor Geometry:\\
Conjectures on Dark Sector, Dark Energy, and Information Geometry}
\author{Marvin L.\ Gentry\\
\small Independent Researcher, Seattle, WA\\
\small \texttt{drmlgentry@protonmail.com}\\
\small ORCID: 0009-0006-4550-2663}
\date{\today}

\begin{document}
\maketitle

\begin{abstract}
We present three conjectures extending the Hyperbolic Flavor Geometry
(HFG) framework beyond the Standard Model flavor sector. The
\emph{Dark Sector Arithmetic Conjecture} predicts that dark matter
particles correspond to non-Lucas primes appearing in covering towers
of cusped hyperbolic 3-manifolds, with specific mass predictions
(e.g., 1.85~MeV sterile neutrino, 39~MeV dark sector mediator,
580~MeV sub-GeV dark matter). The \emph{Cosmological Unfolding
Conjecture} identifies the transition from matter-dominated to
dark-energy-dominated expansion with the spectral transition from
slope-encoded (sparse) to universal geodesics in Dehn fillings of
$m003$ and $m006$, and proposes a structural explanation for the
Hubble tension. The \emph{Ihara--Selberg Bridge} connects the
geodesic length unit $\log\ph$ to the information capacity of a
binary tree, $k\log_2\ph$, suggesting a link between hyperbolic
geometry and error-correcting codes. All conjectures are stated
precisely enough to be falsifiable.
\end{abstract}

\section{Introduction}

The Hyperbolic Flavor Geometry (HFG) framework~\cite{Gentry:CKM,
Gentry:PMNS,Gentry:CP,Gentry:Lucas} provides a geometric origin for
the Standard Model flavor parameters: fermion masses
($m = m_e\ph^{q/4}$), CKM and PMNS mixing matrices, and CP violation
from $\mathbb{Z}/5$ torsion. The framework reduces the flavor sector
to a single arithmetic object: the Lucas sequence
$L_k = \ph^k + \ph^{-k}$ evaluated at quarter-integer indices
$k\in\frac{1}{4}\Z$.

The present note extends the HFG framework to three speculative but
falsifiable directions:
\begin{enumerate}
  \item Dark sector particles identified by non-Lucas primes in
        covering towers.
  \item Dark energy and the Hubble tension as manifestations of
        the spectral transition from sparse to universal geodesics.
  \item An information-geometry correspondence linking the
        geodesic length unit $\log\ph$ to the capacity of binary
        trees and error-correcting codes.
\end{enumerate}

\section{Dark Sector Arithmetic}
\label{sec:dark}

\begin{conjecture}[Dark Sector Arithmetic]
\label{conj:dark}
Let $\mathcal{P}_{\rm visible} = \{2,3,5,7,11,29,47,199,\ldots\}$
be the primes appearing in the covering towers of the visible
sector manifolds $M_{\rm PMNS}$ and $M_{\rm CKM}$ (prime Lucas
numbers together with the ancestor torsion 5).
Any prime $p \notin \mathcal{P}_{\rm visible}$ that appears as
torsion in the covering tower of a cusped hyperbolic 3-manifold
is a candidate for a dark sector particle.

Specifically, the primes $\{13, 17, 19, 23, 31, 41, 43, \ldots\}$
observed in the control scan of $m004$, $m007$, $m009$, $m010$,
$m015$ identify candidate dark sector manifolds.
The mass of a dark fermion associated with such a prime satisfies:
\[
  \frac{m_{\rm dark}}{m_e} = \ph^{q/4},
\]
where $q/4$ is a quarter-integer that is \emph{not} a quarter-Lucas
index (i.e., not in the set $\{k/4 : L_k\in\Z\}$). The covering
tower prime $p$ labels the manifold; the mass index $q$ is a
separate prediction requiring additional structure.
\end{conjecture}

\subsection{Mass predictions for candidate primes}

Using the mass lattice $m = m_e\cdot\ph^{q/4}$, we list the
masses corresponding to the smallest non-Lucas quarter-indices
for each candidate prime, taking $q = p$ as a representative.

\begin{table}[htbp]
\centering
\caption{Dark sector candidate masses from non-Lucas primes.
All masses in MeV unless stated. $\ph = (1+\sqrt{5})/2$.}
\label{tab:dark}
\begin{tabular}{cccl}
\toprule
Prime $p$ & $q = p$ & $m \approx m_e \cdot \ph^{p/4}$ & Note \\
\midrule
13 & 13 & $1.85$ MeV  & sterile neutrino range \\
17 & 17 & $4.8$ MeV   & dark photon range \\
19 & 19 & $6.4$ MeV   & light dark matter \\
31 & 31 & $39$ MeV    & dark sector mediator \\
41 & 41 & $580$ MeV   & sub-GeV dark matter \\
43 & 43 & $940$ MeV   & near proton mass \\
\bottomrule
\end{tabular}
\end{table}

\subsection{Falsifiability: the 17 MeV dark photon anomaly}

The claimed 17~MeV dark photon anomaly (from ${}^8$Be nuclear
transitions) implies mass ratio $m/m_e \approx 33.3$. The
corresponding quarter-index is
\[
  \frac{q}{4} = \frac{\log(33.3)}{\log\ph} \approx 7.28,
  \quad q \approx 29.1.
\]
Since $29 = L_7$ is a prime Lucas number, this anomaly falls in
the \emph{visible} sector in the HFG classification. The
prediction is therefore that if the 17~MeV anomaly represents
a genuine new particle, it is a visible-sector boson, not a
dark sector particle.

\section{Cosmological Unfolding and the Hubble Tension}
\label{sec:cosmo}

\begin{conjecture}[Cosmological Unfolding]
\label{conj:cosmo}
The transition from matter-dominated to dark-energy-dominated
cosmological expansion corresponds to the spectral transition from
the sparse (slope-encoded) tier to the universal tier in the
geodesic length spectra of Dehn fillings of $m003$ and $m006$.

Let $n^*(m003) \approx 21$ and $n^*(m006) \approx 18$ be the
phase transition thresholds (Section~3 of the companion
paper~\cite{Gentry:SU}). The matter-energy equality redshift
$z_{\rm eq} \approx 0.3$ corresponds to these thresholds, which
in turn correspond to Lucas indices $k^* \in \{6.0, 6.4\}$, i.e.,
$L_{k^*} \approx 18$--$23$. The dark energy fraction
$\Omega_\Lambda \approx 0.70$ is numerically close to the average
universal-tier occurrence at threshold:
\[
  \bar{f}(n^*) = \tfrac{1}{2}(65.6\% + 81.8\%) \approx 73.7\%,
\]
a suggestive but not yet derived coincidence.
\end{conjecture}

\subsection{Hubble tension direction}

The lepton manifold $M_{\rm PMNS}$ enters the universal tier at
$n \approx 21$; the quark manifold $M_{\rm CKM}$ enters at
$n \approx 18$. If early-universe (CMB) measurements are dominated
by quark-sector physics and late-universe (distance ladder)
measurements by lepton-sector physics, there is a systematic
offset in the inferred $H_0$.

The observed ratio $H_0^{\rm early}/H_0^{\rm late}
\approx 67.4/73.0 \approx 0.923$ is in the correct direction:
lepton-dominated measurements give lower $H_0$. The
threshold ratio $n^*(m003)/n^*(m006) = 21/18 \approx 1.167$
(lepton threshold higher) is consistent with this direction.

\begin{remark}
The ratio of squared thresholds $(18/21)^2 \approx 0.735$ is
close to $\Omega_\Lambda \approx 0.70$. This numerical coincidence
is recorded for future investigation without claiming a mechanism.
\end{remark}

\subsection{Dehn filling flow model}

A more ambitious formulation models cosmic expansion as a flow
on Dehn filling parameter space:
$(p(t), q(t)) \to (-2,3)$ and $(-5,2)$ as $t \to t_{\rm now}$.
Early universe ($|p|+|q|$ large): near-cusp, universal tier
dominates, symmetric phase, no distinct particles.
Current universe: specific slopes reached, sparse tier emerges,
Standard Model particle content appears. The Hubble parameter
is then
\[
  H(t) \propto \frac{d}{dt}\bigl[\log\mathrm{vol}(M(p(t),q(t)))\bigr],
\]
decreasing as volume approaches the filled manifold values.
Dark energy acceleration corresponds to the slowdown in volume
change as the flow asymptotes to the target slopes.

\section{Ihara--Selberg Bridge and Information Geometry}
\label{sec:ihara}

\begin{conjecture}[Ihara--Selberg Bridge]
\label{conj:ihara}
The geodesic length spectrum of $M_{\rm PMNS}$ is the Riemannian
analog of the Ihara zeta function of a rooted binary tree of
depth $k_{\rm max} = 9$ (the covering tower degree). The
information capacity of a binary tree of depth $k$ is
$C(k) = \log_2 F_k$, where $F_k$ is the $k$-th Fibonacci number.
As $k\to\infty$, $F_k \sim \ph^k/\sqrt{5}$, so
\[
  C(k) \sim k\log_2\ph.
\]
The fundamental geodesic unit $\log\ph$ and the information unit
$\log_2\ph$ are the same object in different bases:
$\log_2\ph = \log\ph / \log 2$.
The prime $2 = L_0$ appears at degree 3 in the covering tower
of $M_{\rm PMNS}$, linking binary information storage to the
arithmetic topology of the manifold.
\end{conjecture}

\begin{remark}
Reed-Solomon error-correcting codes operate over finite fields
$\mathrm{GF}(2^k)$. The base field $\mathrm{GF}(2)$ corresponds
to the prime $2 = L_0$. The covering tower introduces this prime
at degree 3, suggesting a connection between the degree-3 cover
and 3-bit binary encoding. The covering tower of $M_{\rm PMNS}$
may be a geometric model for error-correcting codes with Fibonacci
scaling, though the precise mechanism is not yet understood.
\end{remark}

\section{Discussion}

The three conjectures are speculative but precisely formulated.
The most immediate falsifiable tests are:
\begin{itemize}
  \item Search for dark matter at 1.85~MeV, 6.4~MeV, 39~MeV,
        580~MeV, and 940~MeV; check whether mass ratios to $m_e$
        are Lucas or non-Lucas quarter-powers of $\ph$.
  \item Verify that the 17~MeV anomaly (if real) is a visible-sector
        boson, not a dark sector particle.
  \item Model the Hubble tension using lepton/quark threshold ratio
        $n^*(m003)/n^*(m006) = 21/18$.
  \item Compute the information capacity of the $M_{\rm PMNS}$
        covering tower and compare to known Fibonacci-based
        error-correcting code efficiencies.
\end{itemize}

\section*{Data Availability}
All computations used SnapPy~\cite{SnapPy} version 3.3.2.
Scripts: \url{https://github.com/drmlgentry/hyperbolic-flavor-scan}.

\section*{Conflict of Interest}
The author declares no conflicts of interest.

\begin{thebibliography}{99}

\bibitem{Gentry:CKM}
M.~L.~Gentry,
\textit{Quark Mixing from Hyperbolic Geometry},
Results in Physics, submitted (2026).
SSRN: \href{https://doi.org/10.2139/ssrn.6583550}{10.2139/ssrn.6583550}.

\bibitem{Gentry:PMNS}
M.~L.~Gentry,
\textit{Lepton Mixing from the Borel Structure of Hyperbolic Holonomy},
Results in Physics, submitted (2026).
SSRN: \href{https://doi.org/10.2139/ssrn.6583553}{10.2139/ssrn.6583553}.

\bibitem{Gentry:CP}
M.~L.~Gentry,
\textit{CP Violation from A-Factor Twist Angles of Hyperbolic Holonomy},
Results in Physics, submitted (2026).
SSRN: \href{https://doi.org/10.2139/ssrn.6583551}{10.2139/ssrn.6583551}.

\bibitem{Gentry:Lucas}
M.~L.~Gentry,
\textit{The Lucas Sequence as the Arithmetic Bridge in Hyperbolic
Flavor Geometry},
J.\ Number Theory, submitted (2026). JNTH-D-26-00538.
SSRN: \href{https://doi.org/10.2139/ssrn.6660598}{10.2139/ssrn.6660598}.

\bibitem{Gentry:SU}
M.~L.~Gentry,
\textit{Spectral Universality and Slope-Encoded Arithmetic in Dehn
Fillings of Hyperbolic 3-Manifolds},
in preparation (2026).

\bibitem{SnapPy}
M.~Culler, N.~M.~Dunfield, M.~Goerner, J.~R.~Weeks,
\textit{SnapPy},
\url{http://snappy.computop.org}.

\end{thebibliography}

\end{document}
